% id: 118-02
% book: 베이직쎈 확률과 통계 · 06 이항분포와 정규분포
% section: 유형 074 이항분포와 정규분포의 관계
% figure: none
% answer: 
% answer_source: 
% variant_level: 0 (원본 전사)
% 이 파일은 build.mjs 가 items.json 에서 만든다. 고칠 때는 items.json 을 고친다.
\smash{\raisebox{4.5mm}{\dmkichul{베이직쎈 확률과 통계 118쪽 02번}}}
확률변수 $X$가 이항분포 $\mathrm{B}\left(192{,}\mkern3mu\mathopen{}\dfrac{3}{4}\right)$을 따르면 $X$는 근사적으로 평균이 $m$, 표준편차가 $\sigma$인 정규분포를 따른다. 표준정규분포 $\mathrm{N}(0{,}\mkern3mu\mathopen{}1)$을 따르는 확률변수 $Z$에 대하여\par\smallskip\dispeq{$\mathrm{P}(X\le 132)=\mathrm{P}(Z\le a)$}\par\smallskip\noindent 일 때, $m+\sigma+a$의 값은? \dmcondfill{$a$는 상수이다.}{}
\dmchoicesauto{}{$142$}{$144$}{$146$}{$148$}{$150$}
